\section{Definitions and reconstruction}
\label{app:definitions}

The unit of the main analysis is a security-reconstitution event. An issuer can contribute several observations, and different share classes can have different index status. A review cycle defines the primary inference cluster. AD and ED are dates mapped to the observed trading calendar. ``Latest'' means the 2022--2026 block through the first 2026 review, not a complete five-year realization. The earlier and later decades are 2007--2016 and 2017--2026, respectively.

Returns are daily percentages, and reported CARs sum daily stock returns less OMXSBGI returns. A 2.15\% CAR therefore denotes 2.15 percentage points in the arithmetic cumulative excess-return measure, rather than a claim that a compounded investment outperformed by exactly 2.15\%. Announcement, implementation and primary windows are $[AD-1,AD+1]$, $[ED-1,ED+1]$ and $[AD-1,ED+1]$. The exploratory extension uses $[AD-1,ED-1]$ and $[ED,ED+20]$ on identical complete cases.

For an event-weighted mean $\widehat\mu=N^{-1}\sum_i y_i$, the independent CR1 variance estimator is
\begin{equation}
 \widehat V(\widehat\mu)=\frac{G}{G-1}\frac{1}{N^2}
 \sum_{g=1}^G\left[\sum_{i\in g}(y_i-\widehat\mu)\right]^2.
\end{equation}
The intercept-only regression has one estimated parameter, so the usual $(N-1)/(N-K)$ correction equals one. The joint covariance of the return and ownership means uses the corresponding outer products of the two-dimensional review scores. The Fieller calculation solves the resulting quadratic inequality rather than dividing separate marginal confidence limits.

The cluster jackknife variance is $(G-1)/G$ times the sum of squared differences between the leave-one-review-out estimates and the full estimate. The two-way variance combines separately corrected review and security components and subtracts the corrected review-security intersection. These definitions are supplied to make the calculations inspectable; their validity still depends on the dependence structure and relevant asymptotics.

\section{Weighting and dependence diagnostics}
\label{app:dependence}
\begin{table}[!htbp]
\centering
\caption{Weighting and dependence sensitivity}\label{tab:dependence}
\small
\setlength{\tabcolsep}{5pt}
\begin{tabular}{llrrlrrr}
\toprule
Type & Period & Event & Review & Leave one out & CR1 & CV3 & Two way \\
\midrule
Addition & 2017--2026 & 2.15 & 2.50 & [1.87, 2.47] & 0.55 & 0.55 & 0.47 \\
Addition & 2022--2026 & 1.50 & 1.50 & [1.02, 2.01] & 0.74 & 0.77 & 0.71 \\
Deletion & 2017--2026 & -1.04 & -0.52 & [-1.45, -0.50] & 0.87 & 0.88 & 0.87 \\
Deletion & 2022--2026 & -1.53 & -1.32 & [-2.06, -0.58] & 1.12 & 1.17 & 1.11 \\
\bottomrule
\end{tabular}
\begin{minipage}{\linewidth}\footnotesize\vspace{5pt}
CAR $[AD-1,ED+1]$ in percent. Review mean gives each review equal weight. Leave one out is the range of event-weighted estimates after removing one review at a time, not a confidence interval. CR1, CV3 and Two way report standard errors. CV3 is the cluster jackknife centered on the full-sample estimate. Two way combines review and security clustering with subtraction of the review-security intersection. It does not solve dependence between different securities across different reviews. Source: audit/results/dependence\_sensitivity.csv.
\end{minipage}
\end{table}


The event and review means differ because review sizes vary. Neither weighting scheme is inherently more causal. They answer, respectively, what return the average observed event experienced and what return an average observed review delivered after first averaging its events. The leave-one-out range describes influence, while the standard errors summarize different maintained dependence assumptions. The complete four-block diagnostics are preserved in \texttt{audit/results/dependence\_sensitivity.csv}.
\FloatBarrier

\section{Multiplicity and equivalence}
\label{app:multi}
\begin{table}[!htbp]
\centering
\caption{Multiplicity in a declared ten test audit family}\label{tab:multiplicity}
\small
\setlength{\tabcolsep}{5pt}
\begin{tabular}{llrr}
\toprule
Type & Test & Raw $p$ & Holm $p$ \\
\midrule
Addition & 2007--2011 & 0.027 & 0.214 \\
Addition & 2012--2016 & 0.908 & 1.000 \\
Addition & 2017--2021 & 0.005 & 0.046 \\
Addition & 2022--2026 & 0.062 & 0.310 \\
Addition & 2022--26 minus 2007--11 & 0.029 & 0.214 \\
Deletion & 2007--2011 & 0.007 & 0.059 \\
Deletion & 2012--2016 & 0.522 & 1.000 \\
Deletion & 2017--2021 & 0.759 & 1.000 \\
Deletion & 2022--2026 & 0.279 & 1.000 \\
Deletion & 2022--26 minus 2007--11 & 0.039 & 0.235 \\
\bottomrule
\end{tabular}
\begin{minipage}{\linewidth}\footnotesize\vspace{5pt}
The audit family contains eight five-year total-window means and two latest-minus-earliest contrasts. It was declared after the inherited results were known. Holm adjustment controls familywise error for this family under valid marginal tests; it cannot correct unrecorded historical specification search. It does not turn the thesis into a preregistered study. Source: audit/results/multiplicity.csv.
\end{minipage}
\end{table}


Only the 2017--2021 addition mean remains below 0.05 under the declared ten-test Holm adjustment. The adjustment is conditional on the validity of each underlying marginal test and concerns only this family. It does not include every holdings, mechanism, horizon or historical specification examined during the project. The absence of adjusted significance in other rows neither proves a null effect nor invalidates their point estimates as descriptive summaries.

The latest-period equivalence check uses an audit margin of $\pm1$ percentage point. The analytical 90\% interval is [0.13, 2.87] for additions and [-3.62, 0.55] for deletions. Neither is contained within the margin. The archived approximate MDE for the latest addition mean is 2.35 percentage points and for the later-decade addition mean is 1.62. The last-two-period contrast has an approximate MDE of 3.06 percentage points, and the decade contrast 5.23. These plug-in quantities explain why economically relevant differences can remain unresolved; they are not exact power calculations for the bootstrap or independent evidence of an effect.
\FloatBarrier

\section{Exploratory horizon exclusions}
\label{app:horizon}
\begin{table}[!htbp]
\centering
\caption{Exclusions in the exploratory cumulative return extension}\label{tab:horizon_checks}
\small
\setlength{\tabcolsep}{5pt}
\begin{tabular}{llrrlr}
\toprule
Period & Sample & $N$ & CAR & CR1 interval & $p_{WCR}$ \\
\midrule
2017--2026 & full & 169 & 3.30 & [1.31, 5.28] & $<0.001$ \\
2017--2026 & without Sivers & 168 & 3.07 & [0.96, 5.18] & 0.003 \\
2017--2026 & without MSCI overlap & 157 & 2.51 & [0.60, 4.42] & 0.008 \\
2017--2026 & without screened company news & 163 & 3.14 & [1.13, 5.15] & $<0.001$ \\
2022--2026 & full & 87 & 1.71 & [0.72, 2.69] & 0.010 \\
2022--2026 & without Sivers & 86 & 1.25 & [-0.14, 2.63] & 0.087 \\
2022--2026 & without MSCI overlap & 82 & 1.21 & [-0.25, 2.67] & 0.114 \\
2022--2026 & without screened company news & 84 & 1.79 & [1.04, 2.53] & 0.003 \\
\bottomrule
\end{tabular}
\begin{minipage}{\linewidth}\footnotesize\vspace{5pt}
CAR $[AD-1,ED+20]$ in percent, constructed on complete non-overlapping before/after cohorts. Each exclusion is separate. These follow-up diagnostics were added after inspecting the new cumulative estimates, and inherit the limitations of the original screens. Analytical intervals and 49,999-draw wild-bootstrap $p$-values have the same definitions as Table \ref{tab:horizon}. Source: audit/results/horizon\_sensitivity.csv.
\end{minipage}
\end{table}


The follow-up exclusions are intentionally retained alongside the complete-cohort result. Each changes the population rather than ``correcting'' the baseline into an uniquely valid estimate. In particular, the outcome-informed news screen is incomplete, and MSCI exclusion does not identify the incremental effect of OMXSB membership. The complete output also reports deletion estimates and all four historical blocks, preventing the presentation of the later addition result as the only horizon analyzed.
\FloatBarrier

\section{Auxiliary evidence and its limits}
\label{app:auxiliary}

The issuance exercise relates reconstitution status to the change in shares outstanding from AD-5 to ED+60, after the inherited split adjustment. Its continuous estimate is 0.57 percentage points, with an interval of [-0.47, 1.65] and $p=0.284$. The indicator for an increase exceeding 1\% has an estimated difference of 5.77 percentage points and nominal $p=0.042$; the 5\% threshold gives 3.67 percentage points and $p=0.091$. These thresholds are part of the exploratory record. The adjustment flags 48 of 2,586 rows as split-like, and 23 changes exceeding 30 log percentage points remain. Without reconciliation against all relevant corporate actions, these results do not establish an issuer supply response to index demand. Source: \texttt{other\_tests/issuance\_omxs30\_results.csv} and the archived construction code.

The OMXS30 exercise contains 36 usable events, 19 additions and 17 deletions, across 15 reviews in the pooled sample. The addition mean is 0.21\%, with an interval of [-2.59, 3.27]; the sign-adjusted deletion mean is 0.69\%, with an interval of [-3.35, 7.16]. The archived approximate MDE for the pooled OMXS30-minus-OMXSB contrast is 3.67 percentage points. The results are too imprecise to demonstrate that the indices have equal effects or that either effect is economically zero. They also concern different selection rules and firms.

The turnover-rank exercise compares the characteristics of additions over time. Their absolute average rank changes from approximately 53.0 to 43.1; the inherited contrast has $p=0.261$. Percentile ranks move more because the observed universe expands from roughly 316 to 709 securities. Moreover, the inherited ranking uses six months of turnover near the pre-announcement month, while the current official procedure uses monthly rankings over twelve months. It cannot be interpreted as a validated forecast of the actual Nasdaq selection rule. Source: \texttt{liquidity/liquidity\_results.csv} and the size-rank aggregate files.

The share-class analysis uses sign-adjusted relative log returns, unlike the arithmetic market-adjusted CAR in the main tables. Its event-weighted and firm-weighted summaries should therefore not be mechanically subtracted from the main return means. The placebo and source-change sensitivities are retained in the aggregate twin files, including outcomes that do not support a precise relative-price response.

\section{Reproducibility and unresolved verification}
\label{app:audit}

The manuscript is tied to repository commit \texttt{d2d41666ed21}; the complete identifier is recorded in the accompanying manifest. The original analysis was rerun in a separate directory. Its full runner completed with zero reported execution failures, and the inherited input checks passed for 43 files. An independently written audit reconstructed all three inherited CAR windows for 487 candidates and passed 73 recorded checks. The checks cover sample counts, arithmetic, date mapping, duplicate keys and comparisons of selected estimators. They are neither 73 independent confirmations of the economic hypothesis nor a certification of every source observation.

A comparison of 40 numerical CSV outputs finds 32 matching within the stated $10^{-7}$ tolerance. Eight contain differences in bootstrap intervals or simulated $p$-values. Point estimates, counts and standard errors agree within the audit tolerances. The maximum absolute difference in the main effect table is 0.0043, while some sparse auxiliary composition-ratio intervals differ substantially. One early addition ratio interval changes from [-52.03, 417.95] in the frozen output to [-25.21, 96.37] in the rerun. Another changes from an unbounded interval to a bounded one. Those unstable endpoints are not used to substantiate manuscript conclusions.

The rerun uses Python 3.12.4, NumPy 1.26.4, pandas 2.2.3 and SciPy 1.15.2. These differ from the environment recorded by the inherited project. The bootstrap routines use finite simulation and numerical test inversion; the observed differences are consistent with sensitivity to their numerical implementation, but the audit does not establish a single exclusive cause. The original code and results remain preserved. The new extension, Fieller calculations, checks and generated exhibits are provided separately.

Initial external retrieval recovered 23 of the 39 review notices. Follow-up research opened notices or alternative original attachments for another 12 reviews; three additional reviews were visible only through indexed text, and the 2008 second-half notice remained unresolved. Some accessible documents are secondary republications. Retrieval is separate from complete field-level reconciliation, so these counts do not certify 35 fully verified reviews. A December 2008 correction notice also prints an internally inconsistent effective year; the dataset uses January 2009, but the timing of the initial announcement still requires reconciliation. Original Nasdaq methodology documents and three example thesis PDFs were recovered. Research references were checked against primary publication records or available primary texts, with full-text availability recorded separately.

The accompanying package contains the manuscript, aggregate exhibit inputs, original analysis code, independent audit code, manifests, source links and an issue register. It does not newly redistribute licensed security-level vendor data or third-party full texts. A researcher with lawful access must supply the documented inputs to rerun the complete empirical chain. Checksums identify the version used but are not a substitute for access rights. The authors must still resolve notice discrepancies, reconcile relevant holdings reports and corporate actions, and make the permitted-use and submission decisions themselves.

\section{Writing benchmark and AI disclosure}
\label{app:disclosure}

The organizational benchmark is the BE451 Autumn 2026 syllabus and introductory research material, together with past SSE finance theses available in the school's archive \citep{sse2026}. The syllabus specifies a maximum of 40 single-spaced, 12-point pages excluding appendices; its guidance places the main text at approximately 25--40 pages from introduction to conclusion, including major exhibits, and recommends no table of contents. This draft follows an empirical-paper structure: question and contribution, focused literature, institutional setting and data, methods, results, mechanisms, robustness, discussion and conclusion. This is an application of documented course guidance, not a claim to know the examiner's private preferences or an award committee's future ranking.

The inspected examples include \citet{fager2024}, \citet{holst2023} and \citet{domeij2021}. Their public archive identifiers are 6053, 5734 and 5028. They inform the emphasis on a narrow question, self-contained exhibits and a visible boundary between design and interpretation. No empirical findings or prose are copied from those theses. The accompanying structure note identifies the examples and documented course requirements.

\textbf{AI disclosure for this working document.} OpenAI Codex assisted with source retrieval, computations, code review, exhibits, layout and manuscript composition. The English prose was AI-generated. The course syllabus permits disclosed assistance but prohibits AI directly writing thesis text. This is therefore a research aid, not a compliant submission manuscript. The authors must write the submitted text, verify its argument and sources, disclose assistance accurately and take responsibility for the work.
